\section{Frobenius separation and formal closure}\label{sec:formal}

We first establish a rigidity statement for subgroups of products of a Honda
formal group. A formal equation over the residue field will force such a
subgroup to satisfy a nonzero homomorphism equation. The first step separates
coefficients by taking increasingly deep Frobenius roots.

\subsection{Fields and Honda formal groups}

Fix a prime $p$. Throughout the proof, put
\[
  k=\overline{\Fp},\qquad C=\Cp^{\flat},
\]
and fix the coefficient-field embedding $k\hookrightarrow\OO_C$ inducing
the identification of the residue field with $k$. The field $C$ is
algebraically closed, complete, and of characteristic $p$. We write
$\mm_C=\{c\in C:|c|<1\}$ and use the maximum norm on finite tuples.
Every nonzero element of $k$ has norm one. Consequently a formal power
series over $k$ can be evaluated at any tuple in $\mm_C^r$.

For $s\geq 1$, set $q_s=p^s$ and $k_s=\mathbb F_{q_s}$. The Lubin--Tate formal-module construction for the unramified extension
of $\Qp$ of degree $s$, applied to the series $pT+T^{q_s}$
\cite[Theorem~1]{LT65}, gives on reduction a one-dimensional
commutative Honda formal group $\Gamma_s$ over $k_s$ with
\[
  [p]_{\Gamma_s}(T)=T^{q_s}.
\]
Write
\[
  A_s=\End_k(\Gamma_s),\qquad \Delta_s=A_s[1/p].
\]
The classification of one-dimensional formal groups
\cite[Theorem~18.5.1]{Hazewinkel78} identifies every formal group of
height $s$ over an algebraically closed field of characteristic $p$
with $\Gamma_s$. Dieudonn\'e theory over perfect fields
\cite[Chapter~III, \S8; Chapter~IV, \S\S1,3--4]{Demazure72}
shows that its associated $p$-divisible group is simple up to isogeny
and that $\Delta_s$ is a central division algebra over $\Qp$ of
dimension $s^2$. The isogeny category of $p$-divisible groups over $k$
is semisimple.
These statements include height one and the prime two.

Every endomorphism of $\Gamma_s$ is defined over $k_s$. Indeed, if
$a(T)=\sum_{j\geq1}a_jT^j$ is such an endomorphism, its commutation
with multiplication by $p$ gives
\[
  a(T^{q_s})=a(T)^{q_s},
\]
so $a_j^{q_s}=a_j$ for all $j$. We use formal-group addition in
$\Gamma_s$ and its products; ordinary sums of power-series coefficients
and ordinary differences of coordinates retain their usual meaning.

\subsection{Separating coefficients}

For a multi-index $\ell=(\ell_1,\ldots,\ell_r)\in\mathbb N^r$, put
$|\ell|=\sum_i\ell_i$ and $T^\ell=\prod_i T_i^{\ell_i}$. Roots under
Frobenius in $C$ are unique.
For a related separation argument using high Frobenius substitutions,
see \cite[Proposition~3.1]{Chai08}.

\begin{lemma}[Frobenius separation]\label{formal:separation}
Let $h,r\geq1$, put $q=p^h$, and let $(b_\ell)_{\ell\in\mathbb N^r}$
be an arbitrary family in $C$. Suppose $s\in\mm_C^r$ satisfies, for
every integer $m\geq0$,
\begin{equation}\label{formal:root-bound}
  \limsup_{N\longrightarrow\infty}
  \left|\sum_{|\ell|\leq m} b_\ell^{1/q^N}s^\ell\right|
  \leq \|s\|^{m+1}.
\end{equation}
Then, for every $k$-linear functional $\lambda:C\to k$,
\[
  \sum_{\ell\in\mathbb N^r}\lambda(b_\ell)s^\ell=0.
\]
The series in this conclusion converges, and no continuity assumption
on $\lambda$ is required.
\end{lemma}

\begin{proof}
Fix $m$. If all the coefficients through degree $m$ vanish, the
corresponding assertion is immediate. Otherwise choose a $k$-basis
$\beta_1,\ldots,\beta_a$ for their span and write
\[
  b_\ell=\sum_{j=1}^a c_{\ell j}\beta_j,
  \qquad c_{\ell j}\in k\quad (|\ell|\leq m).
\]
There is a subsequence of positive integers $N$ on which $q^N$-power
Frobenius fixes all the $c_{\ell j}$: finitely many elements of
$\overline{\Fp}$ lie in a common finite field. On this subsequence,
\[
  \sum_{|\ell|\leq m}b_\ell^{1/q^N}s^\ell
  =\sum_{j=1}^a\beta_j^{1/q^N}v_j,
  \qquad v_j=\sum_{|\ell|\leq m}c_{\ell j}s^\ell.
\]
Put $V=\max_j|v_j|$. If $V>0$, choose $j_0$ with $|v_{j_0}|=V$,
set $w_j=v_j/v_{j_0}$, and let $\rho_j\in k$ be the residue of
$w_j\in\OO_C$. At least one $\rho_j$ is nonzero, and
\[
  \eta:=\max_j|w_j-\rho_j|<1.
\]
Pass to a further subsequence on which $q^N$-power Frobenius also fixes
all the $\rho_j$. The independence of the $\beta_j$ gives
$\sum_j\beta_j\rho_j\ne0$, whence
\[
  \left|\sum_j\beta_j^{1/q^N}\rho_j\right|
  =\left|\sum_j\beta_j\rho_j\right|^{1/q^N}
  \longrightarrow1.
\]
On the other hand,
\[
  \limsup_N
  \left|\sum_j\beta_j^{1/q^N}(w_j-\rho_j)\right|
  \leq\eta<1,
\]
since $|\beta_j|^{1/q^N}\to1$ for every $j$. The ultrametric
inequality therefore implies
\[
  \left|\sum_j\beta_j^{1/q^N}v_j\right|\longrightarrow V
\]
along this subsequence. Equation~\eqref{formal:root-bound} gives
$V\leq\|s\|^{m+1}$; this inequality also holds if $V=0$.

For an arbitrary $k$-linear functional $\lambda$, we now have the
finite-sum estimate
\[
  \left|\sum_{|\ell|\leq m}\lambda(b_\ell)s^\ell\right|
  =\left|\sum_j\lambda(\beta_j)v_j\right|
  \leq V\leq\|s\|^{m+1}.
\]
Here $|\lambda(\beta_j)|\leq1$, because these values lie in $k$.
The series with coefficients $\lambda(b_\ell)\in k$ converges at
$s$, and the last bound tends to zero as $m\to\infty$. At $s=0$,
Equation~\eqref{formal:root-bound} for $m=0$ forces $b_0=0$, giving
the same conclusion.
\end{proof}

\begin{remark}\label{formal:separation-quantifiers}
The family $b_\ell$ need not be bounded. All uses of $\lambda$ in
the proof concern finite sums. Moreover, the conclusion holds for
\emph{every} $\lambda$ at each tuple satisfying the hypothesis.
Thus, if the same family satisfies Equation~\eqref{formal:root-bound}
at every point of a set, one may choose a single $\lambda$ nonzero on
one prescribed nonzero $b_\ell$ and obtain a common nonzero formal
equation over $k$ on that set. The Frobenius subsequences used to prove
the estimates may depend on the tuple and on $m$.
\end{remark}

\subsection{Reduced formal subgroups}

To prove formal closure, we will show that the ideal of all equations
vanishing on a subgroup of points defines a reduced formal subgroup.
The next two lemmas turn that formal subgroup into a nonzero
homomorphism equation. We first prove smoothness, before its dimension
or irreducibility is known. The use of ordinary differentials is
justified by the finite Frobenius expansion over a perfect coefficient
field.

\begin{lemma}\label{formal:reduced-smooth}
Let $\mathcal G$ be a smooth formal group over $k$ of dimension $r$,
with coordinate ring $R=k[[T_1,\ldots,T_r]]$. If
$\mathcal H\subseteq\mathcal G$ is a reduced closed formal subgroup,
then $\mathcal H$ is smooth. More precisely, for
$B=R/I$, with augmentation ideal $\mathfrak m_B$, put
$e=\dim_k(\mathfrak m_B/\mathfrak m_B^2)$. Then
$B\cong k[[X_1,\ldots,X_e]]$ as a complete local $k$-algebra.
\end{lemma}

\begin{proof}
Every $f\in R$ has a unique finite decomposition
\[
  f=\sum_{0\leq a_i<p}T_1^{a_1}\cdots T_r^{a_r}f_a(T)^p.
\]
An ordinary $k$-derivation kills all the $p$th powers, so this identity
shows that it is determined by its values on the $T_i$, through the
usual formal partial derivatives. Those partial derivatives themselves
are derivations. Hence the ordinary module $\Omega^1_{R/k}$ is free
on $dT_1,\ldots,dT_r$, and the quotient presentation gives
\begin{equation}\label{formal:differential-presentation}
  \Omega^1_{B/k}
  \cong B^r\big/
  \left\langle
    \bigl(\partial_1f,\ldots,\partial_rf\bigr):f\in I
  \right\rangle_B.
\end{equation}
In particular, this ordinary differential module is finite over $B$.

Consider the completed product ring
\[
  B\widehat\otimes_k B
  =k[[X_1,\ldots,X_r,Y_1,\ldots,Y_r]]/(I(X),I(Y))
\]
and its diagonal ideal $J=(Y_i-X_i)_i$. Writing $Y=X+Z$ modulo
$(Z)^2$ gives
\[
  f(X+Z)=f(X)+\sum_i\partial_i f(X)Z_i\pmod{(Z)^2}.
\]
Thus $J/J^2$ has exactly the presentation in
Equation~\eqref{formal:differential-presentation}.
This computation makes no reducedness assumption on the completed
product. The group-difference automorphism
\[
  (x,y)\longmapsto(x,y-_{\mathcal H}x)
\]
carries the diagonal to the identity section in the second factor.
The conormal module of that section is
$B\otimes_k(\mathfrak m_B/\mathfrak m_B^2)$: it is obtained by
retaining just the linear terms in the second coordinates. We conclude
that
\begin{equation}\label{formal:invariant-differentials}
  \Omega^1_{B/k}\cong
  B\otimes_k(\mathfrak m_B/\mathfrak m_B^2)\cong B^e.
\end{equation}

Let $\mathfrak p$ be a minimal prime of $B$ and set
$D=B/\mathfrak p$, $K=\Frac(D)$, and $d=\dim D$. Since $B$ is
reduced, $B_{\mathfrak p}$ is the field $K$. Localization of ordinary
differentials in Equation~\eqref{formal:invariant-differentials} gives
$\dim_K\Omega^1_{K/k}=e$.

Choose a system of parameters in the complete local domain $D$.
Completeness and the finite-dimensional quotient by those parameters
give a finite map $k[[x_1,\ldots,x_d]]\to D$. It is injective by
dimension, since a nonzero ideal in the source would lower its
dimension. Thus $K$ is a finite extension of
$L=\Frac(k[[x_1,\ldots,x_d]])$. Monomials with exponents below $p$
form a basis over Frobenius, so $[L:L^p]=p^d$. This degree is unchanged
by the finite extension: Frobenius identifies the extensions $K/L$
and $K^p/L^p$, and comparison of the two towers over $L^p$ gives
\[
  [K:K^p]
  =\frac{[K:L]\,[L:L^p]}{[K^p:L^p]}=p^d.
\]
This argument includes inseparable extensions. Because $k$ is perfect,
the $p$-basis description of field differentials now gives
$\dim_K\Omega^1_{K/k}=d$. Indeed, a $p$-basis of $K$ over $K^p$
has $d$ elements, and their differentials form a $K$-basis of
$\Omega^1_{K/k}$. Therefore $d=e$.

Every component has dimension $e$, so $\dim B=e$. This equals the
embedding dimension of $B$, making $B$ regular. Lifting a basis of
$\mathfrak m_B/\mathfrak m_B^2$ gives a surjection
$k[[X_1,\ldots,X_e]]\to B$ by completeness. Equality of dimensions
makes its kernel zero, proving the assertion.
\end{proof}

\begin{lemma}\label{formal:annihilator}
Let $h,r\geq1$. Every proper reduced closed formal subgroup
$\mathcal H\subsetneq\Gamma_h^r$ is annihilated by a nonzero
homomorphism $\Gamma_h^r\to\Gamma_h$ defined over $k$.
\end{lemma}

\begin{proof}
Put $q=p^h$, and let $I\ne0$ be the defining ideal of $\mathcal H$
in $R=k[[T_1,\ldots,T_r]]$. By
Lemma~\ref{formal:reduced-smooth}, its coordinate ring is a power-series
ring $B$ of some dimension $e<r$. If $e=0$, then $B=k$ and
$\mathcal H$ is the origin; a coordinate projection proves the
assertion. Assume henceforth that $e>0$.

Let $\sigma:R\to R$ raise every coefficient to its $q$th power and
fix the variables. For $P\in I$, stability under multiplication by
$p$ gives
\[
  P(T_1^q,\ldots,T_r^q)=(\sigma^{-1}P(T))^q\in I.
\]
Radicality implies $\sigma^{-1}I\subseteq I$. Applying $\sigma$
repeatedly produces the ascending chain
\[
  I\subseteq\sigma I\subseteq\sigma^2 I\subseteq\cdots.
\]
It stabilizes because $R$ is Noetherian, and applying an inverse
iterate of $\sigma$ to a stationary pair gives $I=\sigma I$.
Consequently $\sigma$ descends to an automorphism of $B$.

The coordinate pullback of $[p]$ on $\mathcal H$ is the $h$-fold
absolute Frobenius of $B$, composed with $\sigma^{-1}$. This $h$-fold
Frobenius on $B\cong k[[X_1,\ldots,X_e]]$ is finite free, with basis given by
monomials whose exponents are less than $q$. Thus $[p]$ on
$\mathcal H$ is finite flat of degree $q^e=p^{he}$.
The kernels $\mathcal H[p^j]$ are finite flat of order $p^{hej}$.
The map induced by $[p]$ from level $j+1$ to level $j$ is a pullback
of $[p]:\mathcal H\to\mathcal H$, hence is faithfully flat.
These kernels form a $p$-divisible subgroup $\mathcal H[p^\infty]$
of $\Gamma_h^r[p^\infty]$, of height $he<hr$.

They also recover the formal topology. If $\overline T_i$ are the
images of the ambient coordinates in $B$ and
\[
  J_j=(\overline T_1^{q^j},\ldots,\overline T_r^{q^j}),
\]
then
\begin{equation}\label{formal:kernel-topology}
  \mathfrak m_B^{r(q^j-1)+1}\subseteq J_j
  \subseteq\mathfrak m_B^{q^j}.
\end{equation}
The first inclusion follows because any monomial of the indicated
total degree has an exponent at least $q^j$.

In the semisimple isogeny category over $k$, the ambient group is a
sum of $r$ copies of the simple object associated to $\Gamma_h$.
The subgroup has smaller height. Projecting a nonzero complementary
summand onto one simple factor gives a nonzero rational morphism
\[
  \varphi_{\Qp}:\Gamma_h^r[p^\infty]
  \longrightarrow\Gamma_h[p^\infty]
\]
whose restriction to $\mathcal H[p^\infty]$ is zero.
Choose one common power of $p$ clearing its denominator, and let
$\varphi$ be the resulting actual morphism of $p$-divisible groups.
Its restriction is actually zero. To see this, Hom groups of
$p$-divisible groups are $p$-torsion-free: if $p^a f=0$, then
$f\circ[p]^a=0$, and $[p]^a$ on the source is an epimorphism of
fppf sheaves. A rationally zero restriction is killed by one power
of $p$ in its Hom group and hence vanishes.

The compatible coordinate maps at finite levels recover an actual
formal homomorphism $\Gamma_h^r\to\Gamma_h$. In fact, the ambient
and target kernel ideals are respectively
$(T_1^{q^j},\ldots,T_r^{q^j})$ and $(T^{q^j})$, and inverse limits
of their coordinate rings are the corresponding power-series rings.
The finite-level homomorphism identities pass to these separated
inverse limits. The resulting formal morphism is nonzero since
$\varphi$ is nonzero. Its restriction to $\mathcal H$ vanishes by
Equation~\eqref{formal:kernel-topology} and
$\bigcap_j J_j=0$. This proves the assertion.
\end{proof}

\subsection{From an equation to an annihilator}

\begin{proposition}[Formal closure]\label{formal:closure}
Let $h,r\geq1$, and let $S$ be a subgroup of
$\Gamma_h^r(\mm_C)$. If a nonzero power series in
$k[[T_1,\ldots,T_r]]$ vanishes on $S$, then there are
$a_1,\ldots,a_r\in A_h$, not all zero, such that
\[
  a_1(s_1)+_{\Gamma_h}\cdots+_{\Gamma_h}a_r(s_r)=0
  \qquad\text{for every }s\in S.
\]
\end{proposition}

\begin{proof}
Let $R=k[[T_1,\ldots,T_r]]$, let $F$ be the law of $\Gamma_h^r$,
and let $I\subset R$ be the ideal of all series vanishing on $S$.
Evaluation into the field $C$ shows that $I$ is radical. It is
nonzero by hypothesis and proper because $0\in S$. We show that it
defines a formal subgroup.

For $P\in I$, put $Q(T,T')=P(F(T,T'))$. Fix $s\in S$ and write
\[
  Q(s,T')=\sum_\ell b_\ell(T')^\ell.
\]
Each $b_\ell$ is the evaluation at $s$ of a series over $k$, so
$|b_\ell|\leq1$. For $q=p^h$ and every $t\in S$,
coordinatewise $q^N$th power is multiplication by $p^N$; hence
$Q(s,t^{q^N})=0$. Truncation at degree $m$ gives
\[
  \left|\sum_{|\ell|\leq m}b_\ell t^{q^N\ell}\right|
  \leq\|t\|^{q^N(m+1)}.
\]
Taking $q^N$th roots and applying
Lemma~\ref{formal:separation} shows that, for every $k$-linear
$\lambda:C\to k$, the coefficientwise transform
\[
  Q_{s,\lambda}(T'):=\sum_\ell\lambda(b_\ell)(T')^\ell
\]
vanishes on $S$, and therefore belongs to $I(T')$.

To descend this statement in the first variable, fix $m$ and consider
the finite-dimensional $k$-algebra
\[
  D_m=k[[T']]/\bigl(I(T')+(T')^{m+1}\bigr).
\]
Choose a $k$-basis $e_1,\ldots,e_b$ of $D_m$. Modulo its defining
ideal, write
\[
  Q(T,T')=\sum_{a=1}^b c_a(T)e_a,
  \qquad c_a(T)\in k[[T]].
\]
Every $c_a(T)$ is a finite $k$-linear combination of the coefficients
of $(T')^\ell$ in $Q$, with $|\ell|\leq m$. Consequently
$\lambda(c_a(s))$ is exactly the $e_a$-coordinate of the image of
$Q_{s,\lambda}$ in $D_m$. It is zero. Since algebraic $k$-linear
functionals separate the points of $C$, we obtain $c_a(s)=0$ for
every $s\in S$, and thus $c_a\in I(T)$. This argument moves
$\lambda$ across a finite linear combination of already evaluated
coefficients; it does not move it through an infinite evaluation.

It follows that
\[
  Q\in I(T)+I(T')+(T')^{m+1}
  \qquad\text{for every }m.
\]
Krull intersection \cite[Tag~00IP]{Stacks} in the Noetherian local quotient by
$I(T)+I(T')$ gives $Q\in I(T)+I(T')$. The identity and inversion
conditions follow from $0\in S$ and $-_{\Gamma_h^r}S=S$.
Therefore $\Spf(R/I)$ is a proper reduced closed formal subgroup of
$\Gamma_h^r$.

Lemma~\ref{formal:annihilator} supplies a nonzero formal
homomorphism $\varphi:\Gamma_h^r\to\Gamma_h$ killing this subgroup.
Its restrictions to the factors are endomorphisms
$a_i\in A_h$, and the commutative group law gives
\[
  \varphi(T)=a_1(T_1)+_{\Gamma_h}\cdots+_{\Gamma_h}a_r(T_r).
\]
They are not all zero. The coordinate series of $\varphi$ belongs
to $I$, so evaluating on $S$ proves the required identity.
\end{proof}
